\title{Parameter-Efficient Orthogonal Finetuning via Factorization}

\begin{abstract}
Large foundation models are increasingly prevalent, yet the computational resources required to train them from scratch remain prohibitively high. Consequently, developing strategies for effectively customizing these pre-trained models to specific downstream tasks has become a central research focus. This work examines Orthogonal Finetuning (OFT), a structured and principled approach for adapting models to new tasks. Although OFT demonstrates robust generalization capabilities, it necessitates the optimization of a considerable number of parameters due to the inherently high dimensionality of orthogonal matrices. To tackle this challenge, we begin with an information transmission lens to analyze OFT, subsequently distilling several essential criteria for enhanced parameter efficiency. Drawing inspiration from the efficient data propagation enabled by the Cooley-Tukey fast Fourier transform algorithm, we introduce a novel, compact orthogonal parameterization grounded in butterfly structures. Incorporating this structure into OFT results in Orthogonal Butterfly (BOFT), a new, highly parameter-efficient finetuning method. Notably, BOFT encompasses OFT as a special instance, thereby broadening the orthogonal finetuning paradigm. Extensive experiments are performed to adapt large vision transformers, expansive language models, and state-of-the-art text-to-image diffusion models to a wide range of vision and language downstream tasks.
\end{abstract}

\section{Introduction}

Breakthrough models such as ChatGPT~\cite{brown2020language,chen2021evaluating} and Stable Diffusion~\cite{rombach2022high} have showcased the impressive generalization performance characteristic of large-scale foundation models. This remarkable progress is accompanied by a substantial escalation in parameter count—GPT-3, for example, comprises approximately 175B parameters. The massive scale of these models renders training them from the ground up highly impractical for most researchers. As such, the community’s progress increasingly hinges on the capacity to efficiently repurpose existing foundation models rather than rebuilding them anew. Achieving this goal demands practical approaches for customizing models to downstream tasks without comprehensive retraining. Predominantly, there are three established categories of efficient adaptation techniques: \emph{model finetuning}~\cite{hulora2022,zaken2022bitfit,guo2020parameter,raffel2020exploring,qiu2023controlling,zhang2022adaptive,chavan2023one}, which entails tuning a selected subset of parameters while preserving the original model structure; \emph{adapter tuning}~\cite{rebuffi2017learning,houlsby2019parameter,pfeiffer2020adapterfusion,lin2020exploring,he2021towards}, which augments the base model with trainable modules; and \emph{prompt tuning}~\cite{li2021prefix,lester2021power}, which involves attaching learnable prefix tokens to the input. Among these, model finetuning stands out for its simplicity, efficacy, and—crucially—its maintenance of baseline inference speed.

The core idea underlying model finetuning is to maintain similarity between the pretrained model and its finetuned counterpart according to specific metrics, thereby retaining the knowledge acquired during pretraining. Typically, contemporary finetuning strategies employ a low learning rate, an informal technique intended to minimize the divergence between the two model versions. Recognizing the inherent structure of weight matrices, more systematic approaches seek to conserve the relational characteristics among weight matrices—specifically, the pairwise angular relationships among neurons. Building upon this concept, a new finetuning paradigm termed Orthogonal Finetuning~(OFT)~\cite{qiu2023controlling} has emerged, wherein all neurons within a layer are subject to a transformation via a shared orthogonal matrix, thereby guaranteeing that pairwise neuron angles are preserved throughout the adaptation process. While OFT achieves notable results in terms of generalizability and convergence—particularly when applied to the finetuning of text-to-image diffusion models~\cite{qiu2023controlling}—its reliance on high-dimensional orthogonal matrices results in a large number of parameters to train. To combat this issue, OFT leverages a block-diagonal arrangement for these matrices, reducing parameter count. Nevertheless, such parameter savings introduce trade-offs: the imposed block-diagonal structure enforces a rigid sparsity pattern and necessitates independent orthogonal transformations across different blocks. This constraint may inadvertently result in suboptimal inductive biases, such as diminished representational capacity and the inability to closely approximate standard linear transformations. Furthermore, the challenge of designing or selecting an effective sparsity pattern for these matrices remains unsolved.

The central challenge lies in constructing a dense orthogonal matrix in a manner that remains efficient in terms of parameter usage. At first, this seems daunting, as forming a dense orthogonal matrix of size $d$ traditionally necessitates $\mathcal{O}(d^2)$ parameters. However, our work introduces an alternative by assembling such a matrix through the composition of several factorized sparse matrices. The motivation for this approach is rooted in the principle that parameter count can be minimized by allowing for greater computational complexity—in essence, swapping storage for computational steps. In our method, the orthogonal matrix is not represented directly but rather as a product of sparse matrices, necessitating multiple matrix multiplications in every training epoch. To elucidate this matrix factorization, we frame the construction of dense orthogonal matrices as an exercise in information dissemination. Here, the formation of a dense orthogonal matrix through the multiplication of sparse ones can be likened to the movement of information across a network with grid topology. This analogy provides a conceptual basis for developing various sparse factorization strategies that strike a balance between a reduced number of trainable parameters and the expressive power needed to synthesize dense matrices.

In pursuit of parameter efficiency within this information propagation paradigm, we take inspiration from the butterfly patterns found in the Cooley-Tukey fast Fourier transform algorithm~\cite{cooley1965algorithm}, which are known for their capability to efficiently exchange information among nodes~\cite{klasing1994broadcasting}. The structure of the butterfly graph inherent to the Cooley-Tukey method suggests a sparse matrix factorization technique, termed butterfly factorization. This method allows for the construction of a $d$-dimensional dense matrix through a sequence of $\mathcal{O}(\log d)$ sparse matrices, where each constituent matrix contains merely $\mathcal{O}(d)$ nonzero entries. By imposing orthogonality on every sparse matrix involved, we are able to achieve a highly efficient orthogonal parameterization with only $\mathcal{O}(d\log d)$ parameters—a drastic improvement compared to conventional schemes. Building upon this efficient representation, we introduce Orthogonal Butterfly~(BOFT), a novel approach for parameter-efficient finetuning. BOFT not only generalizes OFT as a special case but also serves as a comprehensive framework for orthogonal finetuning. Notably, both the block-diagonal and butterfly matrix structures share a key property—sparsity—allowing them to lower the number of trainable parameters by introducing structured zeroes. Our strategy stands in interesting contrast to the low-rank formulation employed in LoRA~\cite{hulora2022}; while both low-rank and sparse matrices are cornerstone classes of structured matrices~\cite{candes2011robust} that yield parameter efficiency, they offer distinct pathways to achieving this goal.

In contrast to the block-diagonal architecture employed in OFT to balance model capacity and structural constraints, BOFT leverages the butterfly structure to achieve a more gradual and continuous transition between matrices drawn from the full orthogonal group (for expressivity) and the identity matrix (for regularity). This structural design allows us to restrict the hypothesis space to a more compact subset within the orthogonal group, tailored for specific downstream applications. Given the extensive adoption of the butterfly pattern in efficient linear operations, such as discrete Fourier and cosine transforms, we posit that integrating such structure as a parameter-saving technique inherently biases the model in ways that encourage generalization and mitigate overfitting. This rationale is grounded in the butterfly structure’s ability to replicate a variety of well-known linear transforms—an ability not shared by OFT’s block-diagonal arrangement. Our key contributions are as follows:

\begin{itemize}[leftmargin=*,nosep]
    \item We approach the challenge of parameter-efficient orthogonal finetuning using a novel information transmission perspective. This reimagines the design of parameter-efficient dense orthogonal matrices as an information propagation problem over a grid-graph topology.
    \item Drawing inspiration from the butterfly networks underpinning the Cooley-Tukey algorithm, we introduce Orthogonal Butterfly, a parameter-efficient orthogonal finetuning technique formulated within the information transmission setting.
    \item We present several theoretical arguments elucidating how BOFT achieves a substantial reduction in the number of adjustable parameters while maintaining strong representational power and robust training dynamics. The matrix factorization inherent to BOFT further offers a natural mechanism for weight interpolation (see Figure~\ref{fig:boft_interpolation}).
    \item For the first time, we demonstrate the applicability of orthogonal finetuning in a broad array of domains extending beyond controllable text-to-image synthesis~\cite{qiu2023controlling}, highlighting its viability as a versatile finetuning paradigm. In particular, we deploy BOFT in diverse adaptation tasks across both computer vision and natural language processing. BOFT consistently surpasses prevailing state-of-the-art baselines by a notable margin, thereby establishing its strengths in parameter efficiency and generalization.
\end{itemize}

\section{Related Work}

\textbf{Parameter-efficient finetuning (PEFT)}. With the growing scale and capability of foundation models (\emph{e.g.}, \cite{brown2020language,radford2021learning,rombach2022high,kirillov2023segment}), the pursuit of approaches that allow efficient finetuning for downstream applications has become a prominent topic~\cite{treviso2023efficient,ding2023parameter,lialin2023scaling}. A broad range of PEFT strategies have been proposed~\cite{houlsby2019parameter,aghajanyan2020intrinsic,hulora2022,edalati2022krona,wang2022adamix,gheini2021cross,zaken2022bitfit,guo2020parameter,sung2021training,ansell2022composable,lester2021power,li2021prefix,vu2022spot,he2021towards,mao2021unipelt,karimi2021compacter,liu2022few,sung2022lst,chen2022parameter,jia2022visual,chen2022adaptformer,zhang2022neural,jie2023fact,lian2022scaling,luo2023towards}, of which reparameterization-based techniques~\cite{aghajanyan2020intrinsic,hulora2022,edalati2022krona} are of primary relevance here. For example, LoRA~\cite{hulora2022} enhances pretrained weights by introducing the sum of two low-rank matrices, showing strong results on a range of NLP tasks. In LoRA, a fixed rank is imposed across all inserted matrices; however, alternative works~\cite{zhang2022adaptive,valipour2022dylora,zhang2023increlora} propose adapting ranks per layer for improved allocation of parameter resources. Owing to the straightforward nature of this low-rank reparameterization, it has seen widespread adoption~\cite{dettmers2023qlora,zi2023delta,chavan2023one}. Building on the observation that hyperspherical energy provides insight into generalization~\cite{liu2018learning,liu2021orthogonal}, orthogonal finetuning was recently introduced by~\cite{qiu2023controlling} as an alternative method specifically designed for finetuning text-to-image diffusion models. Here, OFT trains an orthogonal matrix to map activations within a single layer, leading to enhanced generalization and steadier training outcomes relative to LoRA. Nonetheless, OFT typically demands a larger number of trainable parameters when compared to LoRA, highlighting the importance of improving its parameter efficiency. In addition, it remains unexplored whether OFT can be extended to more general adaptation scenarios beyond its original context in text-to-image diffusion models~\cite{qiu2023controlling}. BOFT leverages butterfly factorization to make OFT more parameter-efficient. This development allows us to more broadly evaluate orthogonal finetuning on generic adaptation tasks.

\textbf{Butterfly structures}. The radix-2 Cooley-Tukey method decomposes the $N$-point discrete Fourier transform into two transforms each handling $\frac{N}{2}$ points, naturally forming a butterfly structure, which is representable as a sequence of sparse matrices—a product known as a butterfly matrix. Butterfly matrices have previously been exploited for orthogonal matrix parameterization to bypass pivoting in Gaussian elimination and achieve better computational efficiency~\cite{parker1995random}, as well as for enhancing recurrent neural network training stability~\cite{jing2017tunable} and kernel approximation~\cite{munkhoeva2018quadrature}. Efficient learning of linear transforms via butterfly parameterization has been investigated in~\cite{dao2019learning,dao2019kaleidoscope}, while~\cite{chen2022pixelated,dao2022monarch} employ butterfly matrices for training neural networks more efficiently. Additionally, butterfly structures have been applied to accelerate matrix-vector multiplication~\cite{michielssen1996multilevel,o2010algorithm}, facilitate data-sparse matrix approximations~\cite{li2015butterfly}, and improve network data transmission~\cite{klasing1994broadcasting,li2003linear,rajasingh2016transmission}. Diverging from past applications, our focus is on exploiting butterfly structures to increase the parameter efficiency of OFT.

\section{An Information Transmission View on Orthogonal Finetuning}

We begin by reviewing foundational aspects of OFT. In order to adapt pretrained weight matrices, OFT reformulates the new weights as the product of a trainable orthogonal matrix and the fixed pretrained weight matrix. Whereas LoRA achieves adaptation by adding a low-rank matrix to the existing weights, OFT applies a multiplicative orthogonal matrix for weight modification. While LoRA relies on a low-rank approximation for parameter efficiency, the original OFT implements a block-diagonal orthogonal matrix structure~\cite{qiu2023controlling}; parameter savings increase as the block sizes decrease. Figure~\ref{fig:comp_lora} provides a conceptual illustration of this comparison. The underlying rationale for utilizing orthogonal transformations during finetuning is to maintain the pairwise angles among neurons~\cite{liu2018learning,liu2021orthogonal,qiu2023controlling}, which helps retain the semantic representations learned during pretraining. Specifically, OFT learns an orthogonal matrix $\bm{R}\in\mathbb{R}^{d\times d}$ for a given pretrained linear layer $\bm{W}^0\in\mathbb{R}^{d\times n}$, changing the forward computation from $\bm{z}=(\bm{W}^0)^\top\bm{x}$ to $\bm{z}=(\bm{R}\bm{W}^0)^\top\bm{x}$, where $\bm{x}\in\mathbb{R}^{d}$ and $\bm{z}\in\mathbb{R}^{n}$ denote the input and output vectors, respectively. For zero initialization, OFT sets $\bm{R}$ as the identity matrix at the start. To preserve orthogonality of $\bm{R}$ during training, Cayley parameterization is adopted following \cite{qiu2023controlling,liu2021orthogonal}, that is, $\bm{R}=(\bm{I}+\bm{Q})(\bm{I}-\bm{Q})^{-1}$ where $\bm{Q}$ is a skew-symmetric matrix satisfying $\bm{Q}=-\bm{Q}^\top$. For efficient use of parameters, the orthogonal matrix $\bm{R}$ is given a block-diagonal form: $\text{diag}(\bm{R}_1,\bm{R}_2,\cdots,\bm{R}_r)$, with each $\bm{R}_i\in\mathcal{R}^{b\times b}$ as a compact orthogonal matrix and $br=d$. While the block-diagonal structure boosts parameter efficiency, it rests on the premise that a neuron's dimensions (that is, each column of $\bm{W}^0$) are split into $r$ groups, and each group is independently transformed using its corresponding orthogonal block. Despite observed empirical benefits, this design presupposes an arbitrary partitioning of a neuron's dimensions by their indices, which is not semantically meaningful, thereby highlighting the advantages of dense orthogonal matrices. This prompts the important inquiry: \emph{Is it possible to obtain a dense orthogonal matrix while preserving parameter efficiency?}

In response to this issue, we suggest constructing a dense orthogonal matrix by multiplying several sparse orthogonal matrices. To facilitate a comprehensive yet accessible analysis of the sparsity structure involved in orthogonal matrix factorization, we reinterpret the process of producing a dense orthogonal matrix within OFT as an information transmission scenario. Specifically, if we generate a dense matrix $\bm{R}\in\mathbb{R}^{d\times d}$ as a product of $m$ square matrices, expressed as $\thickmuskip=2mu \medmuskip=2mu \bm{R}=\bm{B}_m\bm{B}_{m-1}\cdots\bm{B}_1$, this sequence can be conceptualized as information propagating across a $d\times (m+1)$ node lattice, as depicted in Figure~\ref{fig:info_trans}. The rationale for this transmission framework stems from the notion that a dense $d$-dimensional square matrix is equivalent to full connectivity between two sets of $d$ nodes. In this context, if a specific entry $R_{ij}$ in $\bm{R}$ equals zero, it signals the impossibility of conveying information from node $j$ to node $i$; conversely, a nonzero $R_{ij}$ enables such communication. Hence, factoring the dense matrix $\bm{R}$ as $\bm{B}_m\bm{B}_{m-1}\cdots\bm{B}_1$ can be seen as orchestrating successive rounds of information transfer based on the graphs encoded by each $\bm{B}_i$. Information is transmitted in the order given by the sequence of matrices, beginning with $\bm{B}_1$ and concluding with $\bm{B}_n$. For instance, Figure~\ref{fig:info_trans} illustrates the factorization $\bm{R}=\bm{B}_5\bm{B}_4\bm{B}_3\bm{B}_2\bm{B}_1$, where both the sparsity configurations and associated graph are shown. The network diagram results from expanding the series of matrix multiplications. Within this graphical model, each $\bm{B}_i$ serves as an adjacency matrix defining connectivity from nodes at the $i$-th level to those at the $(i+1)$-th level. More concretely, the $(j_1, j_2)$ entry in $\bm{B}_i$ indicates the presence or absence (zero indicating absence) of a directed link from the $j_2$-th node of level $i$ to the $j_1$-th node of level $(i+1)$. In order for the matrix $\bm{B}_5\bm{B}_4\bm{B}_3\bm{B}_2\bm{B}_1$ to yield a dense output, it is required that each node at the initial level has the potential to transfer information to every node at the sixth level. If instead we consider only $\bm{R}=\bm{B}_4\bm{B}_3\bm{B}_2\bm{B}_1$ (mapping from the source at level one to the recipient at level five), it becomes evident, for example, that node $1$ at the first level cannot communicate with node $3$ at the fifth level, corresponding to a zero at position $(3,1)$ in the sparsity pattern. A similar alignment between the induced graph and the pattern of non-zeros holds when analyzing products like $\bm{R}=\bm{B}_3\bm{B}_2\bm{B}_1$ or $\bm{R}=\bm{B}_2\bm{B}_1$. In general, for a dense matrix $\bm{R}\in\mathbb{R}^{d\times d}$, the $m$ matrices $\bm{B}_m,\dots,\bm{B}_1$ must collectively define a system of directed edges on a $d\times(m+1)$ grid, with each edge only linking nodes from adjacent columns (levels), ensuring full information transfer from every source at the first level to every node at the $(m+1)$-th level.

The perspective of information transmission offers valuable insight into understanding the sparsity structure observed in the factorization matrices of OFT. As illustrated in Figure~\ref{fig:blk_diag}, the matrix $\bm{R}$ in the conventional OFT exhibits a block-diagonal form. While this arrangement cuts down on the quantity of trainable parameters, it falls short in yielding a densely connected matrix $\bm{R}$. The central aim is to assemble a dense orthogonal matrix using $m$ sparse orthogonal matrices, minimizing the count of effective trainable parameters required. Interpreted through the lens of information transmission, the principal design objectives for this construction are \emph{(i)} \textbf{dense connectivity}, meaning each node from the initial layer must be reachable via some path to every node in the final layer, and \emph{(ii)} \textbf{minimum free edges}, requiring the total number of connections to be kept minimal while preserving orthogonality. Orthogonality, however, imposes intricate restrictions on connections between adjacent layers. In particular, for any matrix $\bm{B}_i$ to achieve full rank—a prerequisite for orthogonality—there must exist $d$ connections to establish a bijection between nodes in the $i$-th and $(i+1)$-th layers, resulting in at least $d$ links per layer transition (for instance, as depicted with 4 edges in Figure~\ref{fig:info_trans}). These $d$ connections are dictated by orthogonality and are non-trainable; for a $d\times d$ orthogonal matrix, the $d$ nonzero entries can take only the values $\pm 1$. Since orthogonal matrices require fewer independent parameters than unconstrained matrices, the orthogonality condition introduces further dependencies among edges. For example, in a $2\times 2$ orthogonal matrix, setting the $(1,1)$ entry to zero necessitates the $(2,2)$ entry to also be zero—i.e., removal of one edge can enforce the removal of another. Consequently, for any permissible set of connections, orthogonality may introduce or eliminate additional edges. When focusing on the pattern of non-zero elements within the resulting orthogonal matrix, the information transmission paradigm becomes particularly advantageous in OFT, as the analysis is limited to the nonzero trainable components of $\bm{R}$, irrespective of their actual values.

A straightforward approach to connecting two layers densely requires $\mathcal{O}(d^2)$ edges, corresponding to one dense orthogonal matrix, resulting in $d^2-d$ edges (for instance, when $d=4$, this totals 12 edges). In contrast, the matrix factorization illustrated in Figure~\ref{fig:info_trans} uses only $10$ edges—fewer than the fully dense orthogonal case. Utilizing this approach offers a graphical viewpoint for analyzing the parameter efficiency of OFT, and allows for the systematic design of efficient factorizations. This line of reasoning is inspired by the butterfly graph topology introduced by the Cooley-Tukey algorithm~\cite{cooley1965algorithm}, which achieves dense connections between $d$ source and $d$ target nodes with only $\mathcal{O}(d\log d)$ edges. As shown in Figure~\ref{fig:info_trans}, $10$ edges suffice to densely connect nodes, but the butterfly network reduces this requirement further to just $8$ edges. In what follows, we present how butterfly structures enhance parameter efficiency.

\section{Orthogonal Parameterization by Butterfly Factorization}

Initially developed for the Cooley-Tukey algorithm in the context of the fast Fourier transform, the butterfly structure enables local modifications in the frequency domain to affect the entire spatial domain—an analogy to our information transmission objective where each first-layer node needs to communicate with all nodes in the final layer. Notably, butterfly structures also serve as prominent topologies in computer networks~\cite{solihin2015fundamentals,li2011linear} to enable efficient data transfer. Letting $k \geq 2$ be a power of $2$, we introduce the \emph{butterfly factor} $\bm{B}^F(k)$ as

\begin{equation}\label{eq:but_factor}
\begin{aligned}
    \bm{B}^F(k)=\begin{bmatrix}
    \text{diag}(\bm{d}_1) & \text{diag}(\bm{d}_2) \\
    \text{diag}(\bm{d}_3) & \text{diag}(\bm{d}_4)
    \end{bmatrix}\in\mathbb{R}^{k\times k},
\end{aligned}
\end{equation}

where each $\bm{d}_i\in\mathbb{R}^{\frac{k}{2}}$ for $i=1,2,3,4$ is a vector. Building on this, and assuming $d=2^N$, the $d$-dimensional \emph{butterfly matrix} $\bm{B}(d)\in\mathbb{R}^{d\times d}$ is recursively constructed as

\begin{equation}
\begin{aligned}
    \!\!\bm{B}(d)=\tilde{\bm{B}}(d,d)\!\cdot\!\begin{bmatrix}
    \bm{B}_1(\frac{d}{2})\!\!\!\!\! & \bm{0} \\
    \bm{0}\!\!\!\!\! &\bm{B}_2(\frac{d}{2})
    \end{bmatrix}= \tilde{\bm{B}}(d,d)\tilde{\bm{B}}(d,\frac{d}{2})\cdots\tilde{\bm{B}}(d,2),
\end{aligned}
\end{equation}

Here, $\bm{B}_1(\frac{d}{2})$ and $\bm{B}_2(\frac{d}{2})$ represent two butterfly matrices each of dimension $\frac{d}{2}$. We introduce the concept of a \emph{butterfly component} as $\thickmuskip=2mu \medmuskip=2mu \tilde{\bm{B}}(d,k)=\text{diag}(\bm{B}^F_1(k),\cdots,\bm{B}^F_{d/k}(k))$, forming a block-diagonal matrix with dimensions $d\times d$ and individual blocks of size $k$, where $\bm{B}^F_i(k),\forall i$ are the butterfly factors as presented in Equation~\ref{eq:but_factor}. Equipped with this structure, we proceed to utilize the butterfly matrix as a tool to parametrize orthogonal matrices. This requires each multiplicative term $\tilde{\bm{B}}(d,k)$ within the butterfly matrix $\bm{B}(d)$ to be orthogonal. Our analysis begins with the block-diagonal matrix $\tilde{\bm{B}}(d,2)$ composed of $2\times 2$ blocks; orthogonality of $\tilde{\bm{B}}(d,2)$ can be achieved via the Cayley transform (or equivalently, $2$-dimensional rotation) applied to each individual block, consistent with \cite{liu2021orthogonal,qiu2023controlling}. Every butterfly component’s sparsity pattern is essentially a permutation of the sparsity pattern found in $\tilde{\bm{B}}(d,2)$; thus, all such butterfly components may be parameterized as orthogonal matrices by analogous means. As a result, we can realize a computationally efficient orthogonal parameterization consisting of multiple $2\times 2$ orthogonal matrices. Extending the butterfly structure according to \cite{chen2022pixelated}, we define the block butterfly component $\tilde{\bm{B}}^b(d,k)$ where each entry of $\bm{d}_i$ is generalized to a $b\times b$ matrix. To ensure orthogonality of $\tilde{\bm{B}}^b(d,2)$, we parameterize every $2b\times 2b$ block matrix to be orthogonal. For butterfly components $\tilde{\bm{B}}^b(d,k)$ with $k>2$, the non-zero structure becomes a block-wise permutation of $\tilde{\bm{B}}^b(d,2)$, thereby enabling their orthogonal parameterization in a similar manner. Aggregating these components, the forward propagation step in BOFT can be described as

\begin{equation}
    \begin{aligned}\nonumber
    \bm{z}=\big(\bm{R}(m,b)\cdot\bm{W}^0\big)^\top\bm{x},~~~~\text{s.t.}~\bigg{\{}\bm{R}(m,b)=\prod_{i=1}^m \tilde{\bm{B}}^b_{(d,i)}~~\&~~\underbrace{\big(\tilde{\bm{B}}^b_{(d,j)}\big)^\top\tilde{\bm{B}}^b_{(d,j)}=\tilde{\bm{B}}^b_{(d,j)}\big(\tilde{\bm{B}}^b_{(d,j)}\big)^\top\!=\bm{I}_d}_{\forall j\in [1, m]}\bigg{\}},
    \end{aligned}
\end{equation}

In the above, for notational clarity, $\tilde{\bm{B}}^b(d,2^{m - i+1})$ is abbreviated as $\tilde{\bm{B}}^b_{(d,i)}$, and $\bm{I}_d$ indicates the $d$-dimensional identity matrix. The orthogonal matrix $\bm{R}(m,b)\in\mathbb{R}^{d\times d}$ is defined as a product of several orthogonal butterfly matrices. For ease of reference, we designate BOFT employing $\bm{R}(m,\frac{b}{2})$ as BOFT($m$,$b$), imposing $b\geq 2$. When setting $\thickmuskip=2mu \medmuskip=2mu m=1$, BOFT($1$,$b$) reduces to a block-diagonal OFT~\cite{qiu2023controlling} parameterized by block size $b$. Likewise, BOFT($1$,$d$) becomes the original OFT~\cite{qiu2023controlling} utilizing a full, unconstrained orthogonal matrix. Selecting $\bm{R}$ as the block butterfly matrix $\bm{B}^{\frac{b}{2}}(d)$ with BOFT($\log \frac{2d}{b}$,$b$) recovers a dense orthogonal matrix $\bm{R}$. More generally, BOFT($m$,$b$) contains $\frac{1}{2}(b-1)dm$ learnable parameters that enable adaptation of a $d\times n$ linear transformation. Employing a butterfly structure with $m=\log d$ and $b=2$, BOFT achieves $\mathcal{O}(d\log d)$ parameter complexity. By contrast, the classical OFT based on a fully dense orthogonal matrix has parameter cost $\mathcal{O}(d^2)$, and the block-diagonal OFT, where the number of blocks is $r$, has $\mathcal{O}(bd)$ parameters. As a result, while the classic OFT requires block size $b=d$ to achieve a fully dense orthogonal transformation, BOFT is able to attain this for any block size $b$.

\textbf{Identity initialization for BOFT}. Standard finetuning strategies aim to closely maintain the behavior of the pretrained model; accordingly, LoRA typically initializes low-rank updates at zero. Following a similar rationale, BOFT initializes every butterfly component as an identity matrix (that is, initializing the skew-symmetric matrices in the Cayley transform to zero).

\textbf{Multiplicative Dropout for BOFT}. LoRA~\cite{hulora2022} integrates a Dropout module for its low-rank updates to mitigate overfitting. Conventional Dropout~\cite{srivastava2014dropout} aligns naturally with LoRA’s additive structure, but not with BOFT’s multiplicative formulation. To remedy this, a dedicated multiplicative Dropout mechanism is proposed for BOFT. Since $\bm{R}(m,b)$ consists of $m$ orthogonal butterfly matrices, each easily rearranged into a $2b\times 2b$ block-diagonal orthogonal form, the multiplicative Dropout works by (i) randomly selecting $p_1$ percent of butterfly matrices and $p_2$ percent of diagonal blocks within each selected butterfly, and (ii) substituting these blocks by identity matrices.

\section{Intriguing Insights and Discussions}

\textbf{Expressivity of BOFT}. While the butterfly structure combined with permutations can exactly recover numerous standard fast linear transforms~\cite{dao2019learning,dao2019kaleidoscope} (\emph{e.g.}, fast Fourier transform, Hadamard transform), the extent to which our orthogonal butterfly matrix can approximate an arbitrary orthogonal matrix is still not well established. To address this, we simulate the approximation of a random dense orthogonal matrix~\cite{anderson1987generation} of size $512\times 512$. The findings, summarized in Figure~\ref{fig:para_eff}, reflect averages over 10 random seeds. The y-axis indicates the approximation error, while the x-axis corresponds to the number of trainable parameters available. Each color in the plot represents BOFT variants using identical block sizes, with the leftmost mark indicating the error from BOFT($1$,$b$) (\emph{i.e.}, standard block-diagonal OFT with block size $b$). BOFT demonstrates improved parameter efficiency relative to OFT. For instance, BOFT($9$,$2$) provides higher expressivity than BOFT($1$,$16$) while utilizing significantly fewer parameters. Typically, choosing a smaller $b$ and a larger $m$ leads to more parameter-efficient models; for example, BOFT($6$,$4$) requires fewer parameters yet achieves an approximation error similar to BOFT($2$,$16$). In general, the butterfly matrix is a more structured subclass within the orthogonal group compared to the block-diagonal form, resulting in BOFT being provably more expressive than OFT for equal block sizes.

\begin{theorem}[Expressivity of BOFT]\label{thm:exp_boft}
    BOFT surpasses OFT in expressivity when both share the same block size. Specifically, to approximate every orthogonal matrix of dimension $d$, one may employ a product of butterfly matrices in the form $\bm{B}_{d-1,1}(d)\bm{B}^\top_{d-1,2}(d)\cdots\bm{B}_{1,1}(d)\bm{B}^\top_{1,2}(d)$, where $\bm{B}_{i,j}(d),\forall i,\forall j$ are butterfly matrices.
\end{theorem}

Theorem~\ref{thm:exp_boft} motivates a straightforward extension for BOFT—one may represent the orthogonal matrix as $\thickmuskip=2mu \medmuskip=2mu \bm{R}^G(m_1,b_1,m_2,b_2,l)=\bm{R}_{l,1}(m_1,b_1)\bm{R}_{l,2}^T(m_2,b_2)\cdots\bm{R}_{1,1}(m_1,b_1)\bm{R}_{1,2}^T(m_2,b_2)$, with $\bm{R}_{i,j}^T(m,b)$ denoting the BOFT block. In particular, for $m_1=m_2=\log d$, $b_1=b_2=2$, and $l=d-1$, the form $\bm{R}^G(m_1,b_1,m_2,b_2,l)$ can capture the entire orthogonal group, a construction termed the kaleidoscope hierarchy~\cite{dao2019kaleidoscope}. Nevertheless, it is important to note that higher expressivity does not automatically guarantee superior results during finetuning; even full finetuning, while being universally expressive, often yields suboptimal outcomes. The balance between expressivity and structural regularization remains critical to robust finetuning. BOFT enhances the finetuning parameter space by incorporating structural priors, thereby allowing a more favorable balance between expressivity and regularity.

\textbf{Spectral properties}. Compared to LoRA, orthogonal finetuning is typically more effective in maintaining desirable spectral characteristics, as it exactly retains the spectral norm of the original pretrained weight matrix $\bm{W}^0$. To illustrate this, consider the singular value decomposition $\bm{W}^0=\bm{U}\bm{\Sigma}\bm{V}^\top$, where $\bm{U}$ and $\bm{V}$ denote orthogonal matrices and $\bm{\Sigma}$ is diagonal containing the singular values. When either OFT or BOFT applies an orthogonal transformation $\bm{R}$ on the left, the resulting weights take the form $\bm{R}\bm{U}\bm{\Sigma}\bm{V}^\top$. This operation leaves the maximal singular value—equivalent to the spectral norm of $\bm{W}^0$—unaltered. The preservation of this property has been found to significantly promote both stability during training and improved generalization~\cite{miyato2018spectral,yoshida2017spectral}. Additional mathematical analysis is available in Appendix~\ref{app:sn_property}.

\textbf{Orthogonal finetuning as learning bilinear similarity}. The BOFT approach can be recast as optimizing a bilinear similarity function, specifically $\bm{w}^0_i\bm{R}\bm{x}$, where $\bm{w}^0_i$ refers to the $i$-th column vector (corresponding to a neuron) in the weight matrix $\bm{W}^0$. From this viewpoint, BOFT essentially learns the bilinear similarity parameterized by matrix $\bm{R}$ under strict regularization (namely, $\bm{R}$ is constrained to be orthogonal). This framework naturally relates to concepts in distance metric learning~\cite{xing2002distance} and bilinear forms~\cite{roman2005advanced}.

\textbf{Inductive bias and generalization in BOFT}. The matrix $\bm{R}(m,b)$ in BOFT often represents a structured subclass of the orthogonal group, introducing specific limitations on the expressiveness of the model. This structure leads BOFT to impart an inherent inductive bias. The particular form of inductive bias arising from butterfly factorization is conjectured to be beneficial for generalization, as it mirrors recurrent structures in numerous established linear transformations~\cite{dao2019kaleidoscope}, for example, the discrete Fourier, sine/cosine, and Hadamard transforms. Additionally, the employment of sparse matrix factorizations in BOFT may yield further, implicit forms of inductive bias~\cite{gunasekar2017implicit,li2018algorithmic,liu2019neural}.

\textbf{Relation to butterfly-based sparse training}. Prior studies~\cite{dao2019learning,dao2019kaleidoscope,chen2022pixelated,dao2022monarch} have explored the use of butterfly parameterization in sparse training. These works predominantly focus on representing the weight matrices with butterfly parameterizations and training neural networks entirely from scratch. \cite{dao2022monarch} extends this by examining the finetuning of pretrained weights: specifically, they project weights onto a form of butterfly matrices and subsequently adapt these projected parts for downstream applications. In contrast, BOFT introduces a fundamentally different finetuning approach, in which weights are transformed through the application of layer-wise shared weight matrices.

\input{rebuttal-results}

\section{Concluding Remarks and Limitations}

This work introduces Orthogonal Butterfly, a broadly applicable parameter-efficient finetuning strategy for foundation models leveraging the butterfly structure. The central idea underlying enhanced parameter-efficiency is to express a dense orthogonal matrix as the product of several sparse orthogonal matrices. To facilitate the discovery of suitable matrix decompositions, a graphical information transmission framework is put forward. Utilizing this framework, the butterfly structure emerges as an effective means to realize sparse orthogonal matrix factorization. Through experiments, we illustrate the empirical strengths of BOFT in finetuning models such as large language models, large vision models, and text-to-image generative systems. Our empirical results also establish the advantage of BOFT as a versatile approach to model finetuning.

Notwithstanding its empirical successes, BOFT has limitations. Because the resultant orthogonal matrix in BOFT is constructed by multiplying several orthogonal matrices, training incurs slightly greater runtime costs compared to OFT. The problem of reducing BOFT's training time efficiency remains unsolved. On the positive side, post-finetuning, the orthogonal matrices learned by BOFT can be seamlessly integrated into the pretrained model without introducing extra inference delay. Additionally, it is not yet clear whether the butterfly structure offers the optimal means for information transmission. Our framework for information transmission allows for cross-disciplinary exploration, particularly insights from computer networking, a field that extensively investigates the transmission efficiency of various network topologies. We anticipate that adopting more optimal network configurations could further improve the composition of dense orthogonal matrices.